\begin{tikzpicture}[
  font=\scriptsize,
  induk/.style={draw, rounded corners, fill=praditagreen!15, align=center,
                text width=10.6cm, minimum height=0.9cm, font=\scriptsize\bfseries},
  gap/.style={draw, rounded corners, fill=blue!8, align=center,
              text width=2.1cm, minimum height=1.45cm, font=\scriptsize\bfseries},
  orang/.style={draw, rounded corners, dashed, fill=red!22, align=center,
                text width=10.6cm, minimum height=0.9cm, font=\scriptsize\bfseries},
  panah/.style={-{Stealth[length=2mm]}, thick},
  urai/.style={gray!70, thin}
]
  \node[induk] (tujuan) at (5.875,0)
    {Satu Pasien, Satu Rekam Medis\\
     {\tiny\mdseries pasien dilayani dengan riwayat yang sama di 5 rumah sakit
      dan 18 klinik}};

  \node[gap, fill=orange!12] (bisnis) at (1.3,-2.0)
    {Bisnis\\{\tiny\mdseries G-B01: pendaftaran lintas fasilitas belum ada}};
  \node[gap] (data) at (4.35,-2.0)
    {Data\\{\tiny\mdseries G-D01: satu pasien, banyak nomor rekam medis}};
  \node[gap] (aplikasi) at (7.4,-2.0)
    {Aplikasi\\{\tiny\mdseries G-A04: integrasi masih \textit{point-to-point}}};
  \node[gap] (teknologi) at (10.45,-2.0)
    {Teknologi\\{\tiny\mdseries G-T03, G-T04: akses dan server terpisah}};

  \node[orang] (orang) at (5.875,-4.1)
    {Orang dan proses\\
     {\tiny\mdseries G-B03: peran \textit{data steward}, SOP pendaftaran baru,
      dan pelatihan petugas, yang sering terlewat karena tidak tampak di
      model aplikasi}};

  \foreach \n in {bisnis, data, aplikasi, teknologi} {
    \draw[urai] (tujuan.south) -- (\n.north);
    \draw[urai, dashed] (\n.south) -- (orang.north -| \n.south);
  }

  \draw[panah] (bisnis) -- node[above, font=\tiny] {butuh} (data);
  \draw[panah] (data) -- node[above, font=\tiny] {butuh} (aplikasi);
  \draw[panah] (aplikasi) -- node[above, font=\tiny] {butuh} (teknologi);
\end{tikzpicture}
